\section{Our Approach}
\label{sec:our}

\subsection{Generation of Prior Knowledge via LLM \label{sec:gen}}

\begin{figure}[htb]
\centering
\includegraphics[width=\linewidth]{images/kgen.png}
\caption{Flowchart of the prior knowledge generation process. The LLM is inspired to generate knowledge in this domain by human-written domain-specific text prompts, which should include both normal-related and abnormal-related knowledge. The feature representation of the compressed knowledge is then obtained through text encoding and dictionary learning.}
\label{fig:kgen}
\end{figure}

An essential aspect of weakly supervised video anomaly detection is to incorporate external knowledge into the training process. As shown in Figure \ref{fig:kgen}], we propose a method for generating prior knowledge using an LLM to provide informative textual prompts related to different types of anomalous scenarios. We then utilize a pre-trained language model, such as BART\citep{a-lewis2019bart}, to embed the textual knowledge, followed by dictionary learning for feature purification and compression. The resulting compact and non-redundant feature representation serves as an external knowledge base for the subsequent training process.

\subsubsection{LLM-Based Knowledge Generation}

Given a video anomaly detection dataset, we first leverage GPT-3\citep{a-brown2020language} to generate textual knowledge with task-specific prompts, which provide high-level semantic descriptions conveying relevant information about both normal and anomalous activities in videos. 
% To generate these prompts, we pass an initial query to the GPT3, mentioning the dataset's specific goal aimed at detecting certain types of anomalies.

For example, suppose we have the following query for a dataset targeting "vehicle accidents":

\begin{quote}
"Generate a textual description for a video sequence containing normal traffic flow and a description for a video sequence containing vehicle accidents."
\end{quote}

Then, GPT-3 generates two textual descriptions: one describing normal traffic flow and the other an anomalous sequence featuring vehicle accidents. These generated texts help exploit the knowledge encapsulated in the LLM and, in turn, improve anomaly detection. Formally, for the aforementioned query, it provides prompts $c_{normal}$ and $c_{anomaly}$.

\subsubsection{Knowledge Embedding via Dictionary Learning}
% However, since the generated knowledge may be structurally and semantically loose, this simple embedding makes it difficult to make full use of the valid information.

After obtaining the textual descriptions of the generated knowledge, BART\citep{a-lewis2019bart} is first used to embed the entire generated text and obtain $f_{normal}$ and $f_{anomaly}$. Since directly incorporating the raw embeddings can be computationally expensive and potentially redundant, we process these embeddings using dictionary learning. This compresses high-dimensional embeddings into a concise dictionary of codewords that effectively capture the core semantic content.

Let $e_i \in \mathbb{R}^d$ be the embedding for the textual prompt $i$ ($f_{normal}$ or $f_{anomaly}$) obtained from GPT-3 and BART, where $d$ is the dimensionality. Next, we train a dictionary $D \in \mathbb{R}^{d \times K}$ with $K$ codewords to sparsely represent the knowledge embeddings:

\begin{equation}
\min_{\mathbf{A}} \frac{1}{2} \| \mathbf{e_{i}} - \mathbf{D} \mathbf{A} \|_F^2 + \lambda \| \mathbf{A} \|_1
\end{equation}

where $\mathbf{A}$ is a coding matrix representing the sparse coding of each sample over the dictionary, $\| \cdot \|_F$ denotes the Frobenius norm, $\| \cdot \|_1$ denotes the L1 norm, and $\lambda$ is a regularization parameter controlling the sparsity of the encoding.

\subsection{Maintaining Consistency in Frame Descriptions \label{sec:maintain}}

\begin{figure}[htb]
\centering
\includegraphics[width=\linewidth]{images/maintain.png}
\caption{With additional scene information provided by the object detection model, the LLM corrects incorrect frame descriptions generated by BLIP.}
\label{fig:maintain}
\end{figure}

As shown in Figure \ref{fig:maintain}], we address the issue of inconsistency in textual descriptions extracted from frames within a video by leveraging additional information from an object detection model. The inconsistency arises from the limitations of the image caption model in capturing fine-grained details and context in video frames. To mitigate this problem, we propose an approach that incorporates the output of the object detection model along with the frame descriptions and jointly processes them using an LLM. The LLM's powerful knowledge representation and logical reasoning capabilities enable us to correct and maintain consistency between the textual descriptions and the visual content of the frames.

Let \( F_i \) be a set of frames within a video, and \( T(F_i) \) be the corresponding textual descriptions generated by BLIP\citep{a-li2022blip}. We aim to incorporate the output of YOLOv8\citep{a-reis2023real}, denoted as \( D(F_i) \), which provides more accurate and detailed information about the objects present in the frames. We denote the corrected textual descriptions as \( \hat{T}(F_i) \), which are generated by the LLM as follows:

\begin{equation}
\label{eq:maintain}
\hat{T}(F_i) = \text{LLM}(T(F_i), D(F_i))
\end{equation}

In Equation (\ref{eq:maintain}), \( \text{LLM}(\cdot) \) denotes the LLM function. It takes the original textual descriptions \( T(F_i) \) and the object detection results \( D(F_i) \) as input and outputs the corrected descriptions \( \hat{T}(F_i) \).

\subsection{Cross-Modal Attention with Prior Knowledge}

\begin{figure*}[htb]
\centering
\includegraphics[width=\linewidth]{images/cross.png}
\caption{Overview of the feature processing method. After the features of different modalities are encoded and concatenated, a bidirectional attention module performs deep feature extraction. The processed features are then optimized by the Knowledge Traction module and used for anomaly detection.}
\label{fig:cross}
\end{figure*}

\begin{figure}[htb]
	\begin{minipage}[t]{\linewidth}
		\centering
		\includegraphics[width=\linewidth]{images/attention.png}
		\subcaption*{(a) In the first stage, visual-text attention performs a coarse fusion of features. In the second stage, text-visual attention fine-tunes the fusion of features.}
		\label{fig:detail:attention}
	\end{minipage}
    % \vspace{15pt}
	\begin{minipage}[t]{\linewidth}
		\centering
		\includegraphics[width=\linewidth]{images/fuse.png}
		\subcaption*{(b) Amplify or suppress the different dimensions of the fused features based on prior knowledge, and further optimize the expressiveness of the features.}
		\label{fig:detail:fuse}
	\end{minipage}
\caption{Module implementation details. (a) The bidirectional attention module. (b) The Knowledge Traction module.}	
\label{fig:detail}
\end{figure}

This section details our cross-modal attention mechanism, which leverages prior knowledge obtained from external textual prompts. Its goal is to enable bidirectional interaction between the visual and textual modalities, allowing them to mutually enhance their representations based on the semantic information provided by the normalized textual prompts.

The architecture consists of three branches: the RGB stream, the optical flow stream, and the text stream. These branches process video frames, optical flow data, and the normalized textual prompts, respectively. Figure \ref{fig:cross} provides an overview of the proposed architecture.

The RGB stream and optical flow stream are responsible for extracting spatio-temporal features from the video frames. These streams utilize pre-trained visual models, such as I3D\citep{a-carreira2017quo}, to capture the visual patterns and motion information present in the video sequences. The text stream is encoded by BART\citep{a-lewis2019bart} and used to complement visual information and facilitate integration with prior knowledge in subsequent processes. Next, textual and visual features are fed into our Visual-Text \& Text-Visual Attention Module for deep feature fusion. Subsequently, the fused features are augmented or suppressed in different regions using prior knowledge, which guides further optimization of the feature representation.

\subsubsection{Visual-Text Attention}

As shown in Figure \ref{fig:detail}](a), the visual and textual features first enter the "Visual-Text Attention" module. The visual features are downsampled by a convolution operation and produce the query ($Q_{v}$) after global average pooling and mapping. In contrast, the text features are passed through N independent 1D convolutional layers for feature compression to reduce computational complexity, producing the key ($K_{t}^{i}$) and value ($V_{t}^{i}$). Independent convolutional layers also allow the model to learn from the text features from multiple perspectives, improving generalization. Finally, $K_{t}^{i}$, $V_{t}^{i}$, and $Q_{v}$ undergo the following attention operation:

\begin{align}
F_{t \mid v} &= \text{Concat}_{i\in \left [ 1, N \right ] }\left ( \frac{Q_{v} \times \left ( K_{t}^{i} \right ) ^{T}}{\sqrt{d_{k}} } \cdot V_{t}^{i} \right ) 
\end{align}

Unlike the classical Transformer\citep{a-vaswani2017attention}, we concatenate the output features of all branches and do not use an additional fully connected layer. This coarse attention operation preserves the features of the original modalities to a greater extent and helps the subsequent process perform deeper feature fusion.

\subsubsection{Text-Visual Attention}

The attention mechanism in the "Text-Visual Attention" module is similar to that in the first stage. The difference is that the features obtained in the first stage are used to generate $Q_{t}$, $$Q_{t} = \sum_{i= 1}^{N} \text{Conv1D}(F_{t \mid v}^{i})$$ and the initial visual features are used to generate $K_{v}$ and $V_{v}$. Since only a coarse feature fusion is performed in the first stage, the visual elements of its output features occupy only a small portion. Therefore, $Q_{t}$ can be considered as a representation of textual features in the second stage of the attention operation. Meanwhile, since it also fuses visual features, the second fusion stage is easier to converge and helps align visual and textual features.

In Figure \ref{fig:detail}](a), the merged feature $F_{v \mid t} =Attention(Q_{t},K_{v},V_{v})$ is reshaped and upsampled to restore the original dimensionality. Then, we design a shuffling method to shift different regions of different channels in the feature. We believe this enhances the model's ability to learn temporal and spatial information. Finally, a convolution operation produces $F_{fused}$.
\subsubsection{Knowledge Traction Feature Fusion}

After deep feature fusion, the resulting features are guided and optimized using prior knowledge, as shown in Figure \ref{fig:detail}](b). First, the feature vectors of normal and anomalous prior knowledge, $K_{n}$ and $K_{a}$, pass through a shared FFN, $FFN_{share}$, to extract common features. The resulting vectors are then summed with $GAP(F_{fused})$ and fed into $FNN_{n}$ and $FNN_{a}$ to extract discrepant features, producing $KV_{n}$ and $KV_{a}$. This process can be described by the following equation: $$KV_{n \parallel a } = FNN_{n \parallel a}(FNN_{share}(K_{n \parallel a}) + GAP(F_{fused}))$$ Then, a sigmoid function is applied to $KV_{n}$ and $KV_{a}$ to represent the degree of adaptation of $F_{fused}$ to $K_{n}$ and $K_{a}$. $F_{fused}$ is scaled according to this degree of adaptation to obtain $F_{n}^{'}$ guided by normal prior knowledge and $F_{a}^{'}$ guided by anomalous prior knowledge, respectively. Finally, the features with high activation in different time domains are selected from $F_{n}^{'}$ and $F_{a}^{'}$ and concatenated to obtain $F_{final}$.

\begin{align}
\label{eq:ktff:a}
F_{n \parallel a}^{'} &= \sigma (KV_{n \parallel a}) \times F_{fused} \\ 
\label{eq:ktff:b}
F_{final} &= \exists_{>\lambda }(F_{n}^{'}) \Join  \exists_{>\beta }(F_{a}^{'}) 
\end{align}

In Equations \ref{eq:ktff:a} and \ref{eq:ktff:b}, $\sigma$ represents the sigmoid function. $\exists_{>\chi}(\cdot)$ denotes selecting the part of the features whose intensity is greater than $\chi$. $\Join$ denotes the concatenation of the selected features while preserving their original order.

\subsection{Optimization Process}

Our optimization objective is adapted from RTFM\citep{a-tian2021weakly} with modifications for our method and consists of three parts: $\mathcal{L}_{cls}$, $\mathcal{L}_{sep}$, and $\mathcal{L}_{limit}$. $\mathcal{L}_{cls}$ denotes the snippet-level classification loss, which can be expressed as $$\mathcal{L}_{\text{cls}}(p = f_{\theta}(F_{final}), y) = \alpha BCE(\left \{ p  \right \}_{k}, y) + BCE(\hat{p}, y)$$, where $f_{\theta}$ denotes the classifier, $\{ \cdot \}_{k}$ denotes the $k$ highest magnitudes in the feature vector, $k=5$, and $\hat{p}$ denotes the remaining features.

$\mathcal{L}_{sep}$ is responsible for the separability of normal and abnormal samples. The distance between features is defined as $D_{k}(f^{+},f^{-}) = \sum_{k} \left \| \left \{ f^{+} \right \}_k  \right \|_{2} - \sum_{k} \left \| \left \{ f^{-} \right \}_k  \right \|_{2}$, and the constraint can be described as $L_{cons}(f^{+}, f^{-}) = \left | m - D(f^{+}, f^{-}) \right | + D_{\lambda \in (+,-)}(\{ f^{\lambda } \}_{k}, \{ f^{\lambda } \}_{k^{'}})$. The final loss is $$\mathcal{L}_{sep} = L_{cons}(F_{t \mid v}^{n}, F_{t \mid v}^{a}) + L_{cons}(F_{v \mid t}^{n}, F_{v \mid t}^{a})$$, where $\|\cdot\|_2$ denotes the L2 norm. $\{ \cdot \}_{k}$ and $\{ \cdot \}_{k^{'}}$ denote the $k$ largest and $k$ smallest magnitudes in the features, respectively.

$\mathcal{L}_{limit}$ constrains the degree to which prior knowledge bootstraps the features, preventing excessive bootstrapping from distorting the features. It is defined as $$\mathcal{L}_{limit} = \sum_{\chi \in (n,a)}^{} \left | \overline{F_{\chi}^{'}} - F_{\chi}^{'}  \right |$$ where $\overline{\cdot}$ denotes the mean of the feature. 

Finally, the overall objective is $\mathcal{L} = \lambda_1 \cdot  \mathcal{L}_{cls} + \lambda_2 \cdot  \mathcal{L}_{sep} + \lambda_3 \cdot  \mathcal{L}_{limit}$. We set $\lambda_1$, $\lambda_2$, and $\lambda_3$ to $1.0$, $0.01$, and $0.005$, respectively.
